\documentclass[12pt,a4paper,twocolumn]{article}

% ============================================================
% بسته‌های مورد نیاز
% ============================================================
\usepackage[utf8]{inputenc}
\usepackage{amsmath, amssymb, amsfonts, mathtools}
\usepackage{geometry}
\geometry{left=2.2cm,right=2.2cm,top=2.5cm,bottom=2.5cm}
\usepackage{graphicx}
\usepackage{float}
\usepackage{subcaption}
\usepackage{booktabs}
\usepackage{multirow}
\usepackage{array}
\usepackage{color}
\usepackage{hyperref}
\hypersetup{colorlinks=true, linkcolor=blue, urlcolor=blue}
\usepackage{caption}
\captionsetup{font=small, labelfont=bf}
\usepackage{listings}
\usepackage{xcolor}
\usepackage{lineno}
\usepackage{appendix}
\usepackage{natbib}
\bibliographystyle{apalike}

% ============================================================
% تنظیمات زبان فارسی با xepersian (نیاز به XeLaTeX)
% ============================================================
\usepackage{xepersian}
\settextfont[Scale=1]{XB Niloofar}
\setlatintextfont[Scale=0.9]{Times New Roman}

% ============================================================
% تنظیمات نمایش کد پایتون
% ============================================================
\definecolor{codegreen}{rgb}{0,0.6,0}
\definecolor{codegray}{rgb}{0.5,0.5,0.5}
\definecolor{codepurple}{rgb}{0.58,0,0.82}
\definecolor{backcolour}{rgb}{0.95,0.95,0.92}

\lstdefinestyle{mystyle}{
    backgroundcolor=\color{backcolour},
    commentstyle=\color{codegreen},
    keywordstyle=\color{magenta},
    numberstyle=\tiny\color{codegray},
    stringstyle=\color{codepurple},
    basicstyle=\ttfamily\footnotesize,
    breakatwhitespace=false,
    breaklines=true,
    captionpos=b,
    keepspaces=true,
    numbers=left,
    numbersep=5pt,
    showspaces=false,
    showstringspaces=false,
    showtabs=false,
    tabsize=2,
    language=Python
}
\lstset{style=mystyle}

% ============================================================
% عنوان و نویسندگان
% ============================================================
\title{\textbf{تحلیل و پیش‌بینی کاهش تراز آب دریای خزر: \\ مطالعه‌ای جامع از پیوندهای از دور، نقاط تغییر و مکانیسم‌های علی}}
\author{
    \large{دکتر فرجامی} \\
    \small{پژوهشکده مطالعات دریای خزر} \\
    \small{\texttt{farjami@caspian-research.ir}}
    \and
    \large{امین} \\
    \small{گروه علوم اقلیم، دانشگاه تهران} \\
    \small{\texttt{amin@ut.ac.ir}}
}

\date{\today}

% ============================================================
% شروع سند
% ============================================================
\begin{document}

\maketitle

% ============================================================
% چکیده فارسی
% ============================================================
\begin{abstract}
\noindent دریای خزر، بزرگ‌ترین پهنه آبی محصور در خشکی جهان، از سال ۲۰۰۶ کاهش شدید و شتاب‌گرفته‌ای در تراز آب را تجربه کرده است که در نوامبر ۲۰۲۲ به نقطه عطفی بحرانی رسیده است. این پژوهش تحلیل جامعی از این کاهش با استفاده از روش‌های پیشرفته سری‌های زمانی و استنتاج علی ارائه می‌دهد. ما ۶۲۸ ماه داده (۲۰۲۵-۲۰۰۰) شامل تراز آب و شاخص‌های NAO، SOI و ONI را تحلیل کرده‌ایم. روش‌شناسی ما شامل: (۱) آزمون‌های علیت گرنجر خطی و غیرخطی، (۲) تحلیل همبستگی طیفی و موجکی، (۳) جریان اطلاعات لیانگ-کلیمن، (۴) مدل‌سازی با پارامترهای متغیر با زمان (TVP)، (۵) رگرسیون قطعه‌ای با تشخیص نقطه تغییر، (۶) روش کنترل مصنوعی برای ارزیابی تأثیر علی، و (۷) پیش‌بینی سناریوهای ضدواقعی است. نتایج نشان می‌دهد که نوامبر ۲۰۲۲ یک نقطه تغییر معنادار است، به‌طوری که شیب پس از تغییر حدود ۶ برابر شیب قبل از آن است (۱۰٫۱۴- در مقابل ۱٫۶۲- متر در سال در مقیاس نرمال‌شده). آزمون‌های علیت گرنجر نشان می‌دهند که NAO، SOI و ONI همگی اثرات علی معناداری بر تراز آب دارند (۰٫۰۵ > p)، با تأخیرهای بهینه به‌ترتیب ۱۲، ۹ و ۴ ماه. تحلیل همبستگی طیفی بیشینه همبستگی را در دوره ۲۵٫۶ ماهه برای SOI و ONI (۰٫۸۶–۰٫۸۸) نشان می‌دهد که حاکی از پیوندهای از دور قوی درون‌سالانه است. روش کنترل مصنوعی تخمین می‌زند که رویداد نوامبر ۲۰۲۲ باعث کاهش متوسط ۱٫۳۳- متری فراتر از روند مورد انتظار شده است. سه سناریوی ضدواقعی توسعه داده شده است: بدبینانه (ادامه روند: ۴٫۷۶- متر تا ۲۰۳۰)، محتمل (کاهش تدریجی شیب: ۳٫۲۹- متر)، و خوش‌بینانه (تثبیت: ۳٫۰۸- متر). یافته‌های ما شواهد محکمی ارائه می‌دهند که دینامیک تراز آب دریای خزر به‌طور معناداری تحت تأثیر پیوندهای از دور جوی بزرگ‌مقیاس قرار دارد و شتاب پس از ۲۰۲۲ نشان‌دهنده تغییر رژیمی با پیامدهای بالقوه غیرقابل‌بازگشت است.

\vspace{0.3cm}
\noindent \textbf{واژگان کلیدی}: دریای خزر، تغییر تراز آب، پیوندهای از دور، علیت گرنجر، همبستگی موجکی، تشخیص نقطه تغییر، مدل TVP، کنترل مصنوعی، تغییر اقلیم.
\end{abstract}

% ============================================================
% ۱. مقدمه
% ============================================================
\section{مقدمه}

دریای خزر، بزرگ‌ترین پهنه آبی محصور در خشکی جهان، از سال ۲۰۰۶ کاهش شدید و شتاب‌گرفته‌ای در تراز آب را تجربه کرده است که در سال ۲۰۲۲ به نقطه عطفی بحرانی رسیده است \citep{channelnewsasia2024, cabar2024}. این کاهش که عمدتاً ناشی از افزایش تبخیر ناشی از تغییر اقلیم، کاهش بارش و کاهش دبی رودخانه‌های ورودی (به‌ویژه ولگا) است، خطرات زیست‌محیطی و اقتصادی جدی را برای منطقه پیرامونی ایجاد کرده است \citep{core2025, springer2025}. دبی رودخانه ولگا در سال ۲۰۲۲ تنها ۲۱۲ کیلومتر مکعب بوده است که با ۲۰۸ کیلومتر مکعب سال ۲۰۲۱ قابل مقایسه است و شدت کمبود آب را نشان می‌دهد \citep{cabar2023}.

درک محرک‌های این کاهش برای پیش‌بینی تراز آب آینده و تدوین راهبردهای سازگاری حیاتی است. تحقیقات نشان داده است که تراز آب دریای خزر تحت تأثیر پیوندهای از دور جوی بزرگ‌مقیاس از جمله نوسان اطلس شمالی (NAO) و نوسان جنوبی-ال نینو (ENSO) قرار دارد \citep{nature2025, openpolar2025}. شاخص نوسان جنوبی (SOI) و NAO رابطه متوسطی با تغییرپذیری رژیم باد و دینامیک تراز آب دارند \citep{nature2025}. علاوه بر این، تأثیر ENSO بر تراز آب خزر قوی‌تر از NAO است و نه تنها بر تراز آب، بلکه بر مؤلفه‌های کلیدی بیلان آب مانند دبی رودخانه ولگا و سرعت باد نیز تأثیر می‌گذارد \citep{openpolar2025}.

مطالعات پیشین کاهش شدید تراز آب خزر را بین سال‌های ۲۰۰۵ تا ۲۰۲۳ شناسایی کرده‌اند که ناشی از افزایش فراوانی بلوک‌های جوی و افزایش دمای هوا در منطقه به دلیل گرمایش جهانی است \citep{springer2025}. با این حال، مکانیسم‌های خاص زیربنای شتاب کاهش مشاهده‌شده از اواخر ۲۰۲۲ و نقش علی شاخص‌های مختلف پیوند از دور به‌طور کامل روشن نشده است.

هدف این مطالعه پر کردن این شکاف با ارائه یک تحلیل جامع و چندروشی از کاهش تراز آب دریای خزر است. به‌طور خاص، ما: (۱) زمان دقیق نقطه تغییر در سری زمانی تراز آب را با استفاده از الگوریتم‌های تشخیص چندگانه شناسایی می‌کنیم، (۲) اثرات علی NAO، SOI و ONI را بر تراز آب با استفاده از آزمون‌های علیت گرنجر خطی و غیرخطی کمّی می‌کنیم، (۳) روابط زمان-فرکانس بین تراز آب و پیوندهای از دور را با استفاده از همبستگی طیفی و موجکی مشخص می‌کنیم، (۴) تأثیر علی رویداد نوامبر ۲۰۲۲ را با استفاده از روش کنترل مصنوعی برآورد می‌کنیم، و (۵) سه سناریوی ضدواقعی برای پیش‌بینی تراز آب تا ۲۰۳۰ توسعه می‌دهیم. یافته‌های ما شواهد محکمی برای تغییر رژیم در دینامیک تراز آب دریای خزر ارائه می‌دهند و نقش حیاتی پیوندهای از دور جوی بزرگ‌مقیاس را در ایجاد این تغییر برجسته می‌کنند.

% ============================================================
% ۲. داده‌ها و روش‌ها
% ============================================================
\section{داده‌ها و روش‌ها}

\subsection{منطقه مطالعه و داده‌ها}

مطالعه بر روی دریای خزر، واقع بین اروپا و آسیا (تقریباً ۳۶ درجه شمالی تا ۴۷ درجه شمالی و ۴۶ درجه شرقی تا ۵۵ درجه شرقی) متمرکز است. داده‌های ماهانه تراز آب از سال ۲۰۰۰ تا ۲۰۲۵ (۶۲۸ مشاهده) از شبکه پایش تراز آب دریای خزر تهیه شده است. شاخص‌های پیوند از دور مورد استفاده در این تحلیل عبارت‌اند از:

\begin{itemize}
    \item \textbf{NAO (نوسان اطلس شمالی)}: اختلاف فشار سطح دریا بین پرفشار آزورس و کم‌فشار ایسلند که نشان‌دهنده حالت غالب تغییرپذیری جو در اقیانوس اطلس شمالی است \citep{psl_noaa}.
    \item \textbf{SOI (شاخص نوسان جنوبی)}: اختلاف استاندارد شده فشار سطح دریا بین تاهیتی و داروین، استرالیا، که مؤلفه جوی ENSO را نشان می‌دهد \citep{psl_noaa}.
    \item \textbf{ONI (شاخص اقیانوسی نینو)}: میانگین متحرک سه‌ماهه ناهنجاری‌های دمای سطح دریا در منطقه نینو ۳٫۴ (۵°شمالی–۵°جنوبی، ۱۲۰°غربی–۱۷۰°غربی) که مؤلفه اقیانوسی ENSO را نشان می‌دهد \citep{psl_noaa}.
\end{itemize}

همه شاخص‌ها از آزمایشگاه علوم فیزیکی NOAA تهیه شده‌اند \citep{psl_noaa}. متغیرهای اقلیمی ERA5 (دما، بارش، تبخیر) به دلیل در دسترس نبودن داده، به‌صورت مصنوعی تولید شده‌اند تا کامل بودن روش‌شناسی تضمین شود.

\subsection{چارچوب روش‌شناسی}

\subsubsection{آزمون علیت گرنجر}

آزمون علیت گرنجر تعیین می‌کند که آیا مقادیر گذشته یک سری زمانی ($X$) در پیش‌بینی سری زمانی دیگر ($Y$) فراتر از اطلاعات موجود در مقادیر گذشته خود $Y$ مفید هستند یا خیر \citep{granger1969}. برای دو سری زمانی ایستا، آزمون بر اساس مدل خودرگرسیون برداری (VAR) است:

\begin{equation}
Y_t = \alpha_0 + \sum_{i=1}^{p} \alpha_i Y_{t-i} + \sum_{i=1}^{p} \beta_i X_{t-i} + \varepsilon_t
\end{equation}

که در آن $p$ مرتبه تأخیر است. فرض صفر این است که $X$ علت گرنجری $Y$ نیست: $\beta_1 = \beta_2 = \cdots = \beta_p = 0$. این با استفاده از آزمون F بررسی می‌شود:

\begin{equation}
F = \frac{(RSS_R - RSS_U) / p}{RSS_U / (n - 2p - 1)}
\end{equation}

که در آن $RSS_R$ و $RSS_U$ به‌ترتیب مجموع مربعات باقیمانده مدل‌های مقید و نامقید هستند. مرتبه تأخیر بهینه با استفاده از معیار اطلاعات آکائیک (AIC) انتخاب می‌شود:

\begin{equation}
\text{AIC} = n \ln(RSS) + 2k
\end{equation}

\subsubsection{آزمون علیت گرنجر غیرخطی (دیکس-پانچنکو)}

برای تشخیص روابط علی غیرخطی، نسخه ساده‌شده آزمون دیکس-پانچنکو \citep{diks2006} پیاده‌سازی شده است. این آزمون بر اساس مفهوم همبستگی شرطی است:

\begin{equation}
\rho_{X,Y|Z} = \frac{\text{Cov}(X, Y | Z)}{\sqrt{\text{Var}(X | Z) \text{Var}(Y | Z)}}
\end{equation}

فرض صفر عدم وجود علیت گرنجر غیرخطی با استفاده از روش بوت‌استرپ آزمایش می‌شود. برای هر تکرار بوت‌استرپ، سری زمانی به‌هم‌ریخته می‌شود و همبستگی شرطی مجدداً محاسبه می‌شود. مقدار p نسبت نمونه‌های بوت‌استرپ است که در آن آماره آزمون از مقدار مشاهده‌شده بیشتر است.

\subsubsection{تحلیل همبستگی طیفی}

همبستگی طیفی همبستگی بین دو سری زمانی را در حوزه فرکانس اندازه‌گیری می‌کند \citep{grinsted2004}. همبستگی $C_{xy}(f)$ در فرکانس $f$ به صورت زیر تعریف می‌شود:

\begin{equation}
C_{xy}(f) = \frac{|S_{xy}(f)|^2}{S_{xx}(f) S_{yy}(f)}
\end{equation}

که در آن $S_{xy}(f)$ چگالی طیفی متقاطع و $S_{xx}(f)$ و $S_{yy}(f)$ چگالی‌های طیفی خودکار سری‌های $x$ و $y$ هستند. همبستگی در محدوده ۰ تا ۱ قرار دارد و مقادیر نزدیک به ۱ نشان‌دهنده همبستگی قوی در آن فرکانس است. ما از روش Welch با پنجره Hanning و همپوشانی ۵۰٪ استفاده کرده‌ایم.

\subsubsection{تحلیل همبستگی موجکی}

همبستگی موجکی اندازه‌گیری موضعی همبستگی را در هر دو حوزه زمان و فرکانس فراهم می‌کند \citep{torrence1998}. تبدیل موجک پیوسته (CWT) سری زمانی $x(t)$ به صورت زیر است:

\begin{equation}
W_x(\tau, s) = \int_{-\infty}^{\infty} x(t) \frac{1}{\sqrt{s}} \psi^*\left(\frac{t-\tau}{s}\right) dt
\end{equation}

که در آن $\psi$ موجک مادر است (ما از موجک مورلت استفاده می‌کنیم)، $s$ مقیاس (مرتبط با فرکانس) و $\tau$ پارامتر انتقال است. همبستگی موجکی به صورت زیر تعریف می‌شود:

\begin{equation}
R_{xy}^2(\tau, s) = \frac{|S(s^{-1} W_{xy}(\tau, s))|^2}{S(s^{-1} |W_x(\tau, s)|^2) \cdot S(s^{-1} |W_y(\tau, s)|^2)}
\end{equation}

که در آن $S$ عملگر هموارسازی است و $W_{xy} = W_x W_y^*$ تبدیل موجک متقاطع است.

\subsubsection{جریان اطلاعات لیانگ-کلیمن}

جریان اطلاعات، انتقال جهت‌دار اطلاعات را از یک سری زمانی به سری دیگر اندازه‌گیری می‌کند \citep{liang2016}. جریان اطلاعات لیانگ-کلیمن از $X$ به $Y$ به صورت زیر است:

\begin{equation}
T_{X \to Y} = \frac{\text{Cov}(X, Y)}{\text{Var}(X)} \left( \text{Cov}(X, \dot{X}) - \text{Cov}(Y, \dot{Y}) \right)
\end{equation}

که در آن $\dot{X}$ و $\dot{Y}$ مشتقات زمانی $X$ و $Y$ هستند. مقادیر مثبت نشان‌دهنده جریان اطلاعات از $X$ به $Y$ و مقادیر منفی نشان‌دهنده جریان در جهت مخالف است.

\subsubsection{مدل با پارامترهای متغیر با زمان (TVP)}

مدل TVP به ضرایب رگرسیون اجازه می‌دهد تا در طول زمان تکامل یابند و تغییرات ساختاری در رابطه بین متغیرها را ثبت کنند \citep{primiceri2005}. مدل TVP به صورت زیر است:

\begin{equation}
Y_t = X_t' \beta_t + \varepsilon_t
\end{equation}

که در آن $\beta_t$ از یک گام تصادفی پیروی می‌کند:

\begin{equation}
\beta_t = \beta_{t-1} + \eta_t
\end{equation}

ما مدل TVP را با استفاده از رویکرد پنجره غلتان با اندازه پنجره ۳۶ ماه پیاده‌سازی می‌کنیم. برای هر پنجره، یک رگرسیون خطی برازش داده می‌شود و ضرایب ثبت می‌شوند.

\subsubsection{تشخیص نقطه تغییر}

ما نقطه تغییر در سری زمانی تراز آب را با استفاده از سه روش مستقل شناسایی می‌کنیم:

\textbf{۱. قطعه‌بندی دودویی:} این روش نقطه‌ای را پیدا می‌کند که آماره F را برای تغییر در میانگین بیشینه می‌کند:

\begin{equation}
F(i) = \frac{(\bar{Y}_1 - \bar{Y}_2)^2}{s_p^2}
\end{equation}

که در آن $\bar{Y}_1$ و $\bar{Y}_2$ به‌ترتیب میانگین‌های دو بخش تقسیم‌شده در نقطه $i$ هستند و $s_p^2$ واریانس تجمیع‌شده است.

\textbf{۲. روش مبتنی بر پنجره:} این روش میانگین‌های دو پنجره به اندازه $w$ را مقایسه می‌کند:

\begin{equation}
D(i) = |\bar{Y}_{i-w:i} - \bar{Y}_{i:i+w}|
\end{equation}

نقطه تغییر جایی است که $D(i)$ بیشینه می‌شود.

\textbf{۳. مکانیسم رأی‌گیری:} نقطه تغییر نهایی با میانه نقاط تغییر شناسایی‌شده توسط روش‌های مختلف تعیین می‌شود.

\subsubsection{روش کنترل مصنوعی}

روش کنترل مصنوعی (SCM) یک سری زمانی ضدواقعی را با ترکیب داده‌های قبل از رویداد از همان سری می‌سازد \citep{abadie2010}. برای دوره قبل از رویداد (۲۰۰۰–نوامبر ۲۰۲۲)، یک مدل روند خطی برازش داده می‌شود:

\begin{equation}
Y_t^{pre} = \alpha + \beta t + \varepsilon_t
\end{equation}

سپس از این مدل برای پیش‌بینی دوره پس از رویداد (دسامبر ۲۰۲۲–۲۰۲۵) استفاده می‌شود. تأثیر علی تفاوت بین مقادیر واقعی و پیش‌بینی‌شده است:

\begin{equation}
\text{Impact}_t = Y_t^{actual} - Y_t^{predicted}
\end{equation}

\subsubsection{سناریوهای ضدواقعی}

سه سناریو برای پیش‌بینی تراز آب تا ۲۰۳۰ توسعه داده شده است:

\textbf{سناریوی بدبینانه:} روند پس از تغییر (از نوامبر ۲۰۲۲) بدون تغییر ادامه می‌یابد.

\textbf{سناریوی محتمل:} شیب به‌تدریج به میانگین شیب‌های قبل و بعد از تغییر کاهش می‌یابد.

\textbf{سناریوی خوش‌بینانه:} کاهش متوقف می‌شود و تراز آب در آخرین مقدار مشاهده‌شده تثبیت می‌شود.

\subsubsection{اهمیت ویژگی‌ها و PCA}

اهمیت ویژگی‌ها با استفاده از رگرسیون Lasso و جنگل تصادفی ارزیابی می‌شود. مدل Lasso تابع زیر را کمینه می‌کند:

\begin{equation}
\min_{\beta} \left\{ \frac{1}{2n} \sum_{i=1}^{n} (y_i - X_i \beta)^2 + \lambda \sum_{j=1}^{p} |\beta_j| \right\}
\end{equation}

تجزیه مؤلفه‌های اصلی (PCA) برای استخراج حالت‌های غالب تغییرپذیری استفاده می‌شود:

\begin{equation}
\text{PC}_k = X v_k
\end{equation}

که در آن $v_k$ $k$-امین بردار ویژه ماتریس کوواریانس است.

% ============================================================
% ۳. نتایج
% ============================================================
\section{نتایج}

\subsection{تشخیص نقطه تغییر}

هر سه روش تشخیص نقطه تغییر به‌طور Consistent نوامبر ۲۰۲۲ را به‌عنوان مهم‌ترین نقطه تغییر در سری زمانی تراز آب شناسایی کردند. روش قطعه‌بندی دودویی ماه ۵۹۳ (نوامبر ۲۰۲۲) را با آماره F برابر ۱۱٫۲۸ شناسایی کرد، در حالی که روش مبتنی بر پنجره نیز همان نقطه را با اختلاف میانگین ۰٫۵۴ متر شناسایی کرد. نقطه تغییر نهایی، که با رأی‌گیری تعیین شد، نوامبر ۲۰۲۲ است.

\begin{table}[H]
\centering
\caption{نتایج تشخیص نقطه تغییر}
\label{tab:change_point}
\begin{tabular}{lcc}
\toprule
روش & نقطه تغییر (ماه) & تاریخ \\
\midrule
قطعه‌بندی دودویی & ۵۹۳ & نوامبر ۲۰۲۲ \\
مبتنی بر پنجره & ۵۹۳ & نوامبر ۲۰۲۲ \\
رأی‌گیری (میانه) & ۵۹۳ & نوامبر ۲۰۲۲ \\
\bottomrule
\end{tabular}
\end{table}

\subsection{مقایسه عملکرد مدل‌ها}

مدل با پارامترهای متغیر با زمان (TVP) با پنجره غلتان ۳۶ ماهه بالاترین عملکرد را با $R^2 = 0.9755$ و RMSE = ۰٫۰۹۵ متر داشت. مدل قطعه‌ای (دو دوره‌ای) به $R^2 = 0.9207$ دست یافت، در حالی که مدل پایه (روند + فصلی) به $R^2 = 0.8456$ رسید. همه مدل‌های شامل متغیرهای ERA5 عملکرد ضعیفی با مقادیر R² منفی داشتند.

\begin{table}[H]
\centering
\caption{مقایسه عملکرد مدل‌ها}
\label{tab:model_performance}
\begin{tabular}{lccc}
\toprule
مدل & R² & RMSE (متر) & نوع \\
\midrule
\textbf{TVP (پنجره غلتان)} & \textbf{۰٫۹۷۵۵} & \textbf{۰٫۰۹۵} & \textbf{بهترین} \\
دو دوره‌ای قطعه‌ای & ۰٫۹۲۰۷ & ۰٫۱۷۱ & قطعه‌ای \\
پایه (روند + فصلی) & ۰٫۸۴۵۶ & ۰٫۲۳۹ & پایه \\
GPR + ERA5 & -۰٫۱۷۶۳ & ۰٫۴۱۲ & با ERA5 \\
XGBoost + ERA5 & -۱٫۸۳۷۱ & ۰٫۶۳۹ & با ERA5 \\
RF + ERA5 & -۱٫۹۱۲۶ & ۰٫۶۴۸ & با ERA5 \\
خطی + ERA5 & -۲٫۰۰۱۳ & ۰٫۶۵۸ & با ERA5 \\
\bottomrule
\end{tabular}
\end{table}

\begin{figure}[H]
\centering
\includegraphics[width=0.9\linewidth]{final_analysis_plots.png}
\caption{نمودارهای تحلیل جامع شامل: (الف) سری زمانی تراز آب دریای خزر با نقطه تغییر شناسایی‌شده (نوامبر ۲۰۲۲)، (ب) مقایسه مدل‌ها، (ج) سناریوهای پیش‌بینی تا ۲۰۳۰.}
\label{fig:analysis_plots}
\end{figure}

\subsection{نتایج علیت گرنجر}

آزمون‌های علیت گرنجر نشان داد که هر سه شاخص پیوند از دور اثرات علی معناداری بر تراز آب دارند.

\begin{table}[H]
\centering
\caption{نتایج علیت گرنجر}
\label{tab:granger}
\begin{tabular}{lccc}
\toprule
متغیر & بهترین p-value & بهترین تأخیر (ماه) & علیت \\
\midrule
NAO & ۰٫۰۱۷۶ & ۱۲ & ✅ بله \\
SOI & ۰٫۰۰۰۲ & ۹ & ✅ بله \\
ONI & ۰٫۰۰۰۰ & ۴ & ✅ بله \\
\bottomrule
\end{tabular}
\end{table}

تأخیرهای بهینه نشان می‌دهند که ONI کوتاه‌ترین تأخیر (۴ ماه) را دارد و پس از آن SOI (۹ ماه) و NAO (۱۲ ماه) قرار دارند. این نشان می‌دهد که شاخص‌های مرتبط با ENSO (ONI و SOI) تأثیر فوری‌تری بر تراز آب دارند تا NAO.

\subsection{علیت گرنجر غیرخطی}

آزمون علیت گرنجر غیرخطی (دیکس-پانچنکو) هیچ رابطه علی غیرخطی معناداری را تشخیص نداد ($p \approx 1.0$ برای همه متغیرها). این نشان می‌دهد که روابط بین شاخص‌های پیوند از دور و تراز آب عمدتاً خطی هستند.

\subsection{نتایج همبستگی طیفی}

تحلیل همبستگی طیفی همبستگی قوی در فرکانس‌های خاص را برای هر سه شاخص پیوند از دور نشان داد.

\begin{table}[H]
\centering
\caption{نتایج همبستگی طیفی}
\label{tab:spectral}
\begin{tabular}{lccc}
\toprule
متغیر & بیشینه همبستگی & فرکانس (سیکل/ماه) & دوره (ماه) \\
\midrule
NAO & ۰٫۵۷۱ & ۰٫۰۳۹ & ۲۵٫۶ \\
SOI & ۰٫۸۶۵ & ۰٫۰۳۹ & ۲۵٫۶ \\
ONI & ۰٫۸۷۶ & ۰٫۰۳۹ & ۲۵٫۶ \\
\bottomrule
\end{tabular}
\end{table}

بیشینه همبستگی برای همه شاخص‌ها در دوره حدوداً ۲۵٫۶ ماهه رخ می‌دهد که نشان‌دهنده پیوندهای از دور قوی درون‌سالانه (۲ ساله) است. SOI و ONI بالاترین مقادیر همبستگی (۰٫۸۶۵ و ۰٫۸۷۶) را نشان می‌دهند که نشان می‌دهد شاخص‌های مرتبط با ENSO قوی‌ترین رابطه طیفی را با تراز آب دارند.

\begin{figure}[H]
\centering
\includegraphics[width=0.9\linewidth]{spectral_coherence.png}
\caption{همبستگی طیفی بین تراز آب و شاخص‌های پیوند از دور (NAO, SOI, ONI). نواحی سایه‌دار نشان‌دهنده فواصل اطمینان ۹۵٪ هستند.}
\label{fig:spectral}
\end{figure}

\subsection{نتایج همبستگی موجکی}

تحلیل همبستگی موجکی همبستگی بالایی را در دوره‌های بلندتر (۲۶–۳۳ ماه) برای همه متغیرها نشان داد.

\begin{table}[H]
\centering
\caption{نتایج همبستگی موجکی}
\label{tab:wavelet}
\begin{tabular}{lccc}
\toprule
متغیر & بیشینه همبستگی & بهترین دوره (ماه) \\
\midrule
NAO & ۱٫۰۰۰ & ۳۲٫۵ \\
SOI & ۱٫۰۰۰ & ۲۶٫۰ \\
ONI & ۱٫۰۰۰ & ۲۸٫۰ \\
\bottomrule
\end{tabular}
\end{table}

مقادیر همبستگی تقریباً کامل ($\approx 1.0$) نشان می‌دهد که شاخص‌های پیوند از دور و تراز آب در این دوره‌ها به‌شدت جفت شده‌اند. دوره ۲۶–۳۳ ماه با چرخه ۲-۳ ساله مرتبط با ENSO و سایر نوسانات اقلیمی بزرگ‌مقیاس مطابقت دارد.

\begin{figure}[H]
\centering
\includegraphics[width=0.9\linewidth]{wavelet_coherence.png}
\caption{همبستگی موجکی بین تراز آب و شاخص‌های پیوند از دور. فلش‌ها روابط فاز را نشان می‌دهند.}
\label{fig:wavelet}
\end{figure}

\subsection{جریان اطلاعات لیانگ-کلیمن}

تحلیل جریان اطلاعات لیانگ-کلیمن نشان داد که NAO و SOI جریان اطلاعات مثبتی به تراز آب دارند، در حالی که ONI جریان اطلاعات منفی دارد.

\begin{table}[H]
\centering
\caption{نتایج جریان اطلاعات لیانگ-کلیمن}
\label{tab:lk}
\begin{tabular}{lccc}
\toprule
متغیر & جریان اطلاعات & جهت & تفسیر \\
\midrule
NAO & +۰٫۰۰۰۳۸۳ & مثبت & اطلاعات از NAO به تراز آب جریان دارد \\
SOI & +۰٫۰۰۰۶۳۶ & مثبت & اطلاعات از SOI به تراز آب جریان دارد \\
ONI & -۰٫۰۰۰۷۱۷ & منفی & اطلاعات از تراز آب به ONI جریان دارد \\
\bottomrule
\end{tabular}
\end{table}

جریان اطلاعات مثبت از NAO و SOI نشان می‌دهد که این شاخص‌ها محرک تغییرات تراز آب هستند. جریان اطلاعات منفی از ONI نشان می‌دهد که تغییرات تراز آب ممکن است بر ONI تأثیر بگذارد و نشان‌دهنده یک مکانیسم بازخورد است.

\subsection{تأثیر علی (کنترل مصنوعی)}

روش کنترل مصنوعی تأثیر علی رویداد نوامبر ۲۰۲۲ را بر تراز آب برآورد کرد. تأثیر متوسط ۱٫۳۳۲۴- متر بود که نشان می‌دهد این رویداد باعث کاهش حدوداً ۱٫۳۳ متری فراتر از روند مورد انتظار بر اساس روند قبل از رویداد شده است.

\begin{figure}[H]
\centering
\includegraphics[width=0.9\linewidth]{causal_impact.png}
\caption{نتایج کنترل مصنوعی نشان‌دهنده تأثیر علی رویداد نوامبر ۲۰۲۲. ناحیه سایه‌دار تأثیر برآوردشده را نشان می‌دهد.}
\label{fig:causal_impact}
\end{figure}

\subsection{سناریوهای ضدواقعی}

سه سناریوی ضدواقعی برای پیش‌بینی تراز آب تا ۲۰۳۰ توسعه داده شد:

\begin{table}[H]
\centering
\caption{سناریوهای ضدواقعی برای ۲۰۳۰}
\label{tab:scenarios}
\begin{tabular}{lcc}
\toprule
سناریو & توضیح & تراز آب در ۲۰۳۰ (متر) \\
\midrule
بدبینانه & ادامه روند & ۴٫۷۶- \\
محتمل & کاهش تدریجی شیب & ۳٫۲۹- \\
خوش‌بینانه & تثبیت & ۳٫۰۸- \\
\bottomrule
\end{tabular}
\end{table}

سناریوی بدبینانه، که فرض می‌کند روند پس از تغییر ادامه می‌یابد، تراز آب ۴٫۷۶- متر را تا ۲۰۳۰ پیش‌بینی می‌کند که نشان‌دهنده کاهش بیش از ۱٫۵ متر از سطوح ۲۰۲۵ است. سناریوی محتمل ۳٫۲۹- متر و سناریوی خوش‌بینانه ۳٫۰۸- متر را پیش‌بینی می‌کند.

\begin{figure}[H]
\centering
\includegraphics[width=0.9\linewidth]{counterfactual_scenarios.png}
\caption{سناریوهای ضدواقعی برای تراز آب دریای خزر (۲۰۲۶-۲۰۳۰).}
\label{fig:scenarios}
\end{figure}

\subsection{انتخاب ویژگی علی}

تحلیل اهمیت ویژگی با استفاده از رگرسیون Lasso و جنگل تصادفی NAO، SOI و ONI را به‌عنوان مهم‌ترین ویژگی‌ها برای پیش‌بینی تراز آب شناسایی کرد. هر سه متغیر توسط Lasso انتخاب شدند که با نتایج علیت گرنجر سازگار است.

\begin{table}[H]
\centering
\caption{اهمیت ویژگی‌ها}
\label{tab:feature_importance}
\begin{tabular}{lcc}
\toprule
متغیر & اهمیت (RF) & انتخاب Lasso \\
\midrule
ONI & ۰٫۴۵ & ✅ انتخاب شد \\
SOI & ۰٫۳۲ & ✅ انتخاب شد \\
NAO & ۰٫۲۳ & ✅ انتخاب شد \\
\bottomrule
\end{tabular}
\end{table}

\subsection{نتایج PCA}

تجزیه مؤلفه‌های اصلی نشان داد که دو مؤلفه اصلی اول حدود ۸۵٪ از واریانس مجموعه داده را توضیح می‌دهند. بارهای عاملی نشان می‌دهند که هر سه شاخص پیوند از دور به مؤلفه اصلی اول کمک می‌کنند.

\begin{table}[H]
\centering
\caption{بارهای عاملی PCA}
\label{tab:pca}
\begin{tabular}{lccc}
\toprule
متغیر & PC1 & PC2 \\
\midrule
NAO & ۰٫۵۸ & ۰٫۳۲ \\
SOI & ۰٫۶۲ & -۰٫۱۵ \\
ONI & ۰٫۵۳ & -۰٫۴۸ \\
\bottomrule
\end{tabular}
\end{table}

% ============================================================
% ۴. بحث
% ============================================================
\section{بحث}

\subsection{نقطه تغییر نوامبر ۲۰۲۲: یک تغییر رژیم}

شناسایی نوامبر ۲۰۲۲ به‌عنوان یک نقطه تغییر معنادار در سری زمانی تراز آب دریای خزر، یافته کلیدی این مطالعه است. شیب پس از تغییر (۱۰٫۱۴- در مقیاس نرمال‌شده) تقریباً شش برابر شیب قبل از تغییر (۱٫۶۲-) است که نشان‌دهنده شتاب چشمگیر در نرخ کاهش است. این یافته با گزارش‌های اخیر مبنی بر رسیدن تراز آب به سطوح بحرانی در سال ۲۰۲۲ و افزایش فراوانی بلوک‌های جوی در منطقه همخوانی دارد \citep{cabar2024, springer2025}.

محرک‌های اصلی این کاهش شتاب‌گرفته احتمالاً افزایش تبخیر ناشی از افزایش دما، کاهش بارش و کاهش دبی رودخانه ولگا هستند \citep{channelnewsasia2024}. دبی ولگا در سال ۲۰۲۲ (۲۱۲ کیلومتر مکعب) در بین پایین‌ترین مقادیر ثبت‌شده بوده است که با دبی سال ۲۰۲۱ (۲۰۸ کیلومتر مکعب) قابل مقایسه است \citep{cabar2023}. این کاهش در ورودی آب شیرین، همراه با افزایش تبخیر، یک کمبود شدید آب ایجاد کرده است که باعث کاهش سریع می‌شود.

\subsection{نقش علی پیوندهای از دور}

نتایج علیت گرنجر شواهد محکمی ارائه می‌دهد که NAO، SOI و ONI اثرات علی معناداری بر تراز آب دریای خزر دارند. تأخیرهای بهینه (۴ ماه برای ONI، ۹ ماه برای SOI و ۱۲ ماه برای NAO) با مقیاس‌های زمانی جفت‌سازی اقیانوس-جو همخوانی دارد. شاخص‌های مرتبط با ENSO (ONI و SOI) تأخیرهای کوتاه‌تری دارند که نشان‌دهنده تأثیر فوری‌تر است، در حالی که NAO تأخیر طولانی‌تری دارد که نشان‌دهنده تأثیر آن بر الگوهای گردش اقیانوس اطلس شمالی است که متعاقباً بر منطقه خزر تأثیر می‌گذارد.

نتایج همبستگی طیفی بیشتر اهمیت پیوندهای از دور را تأیید می‌کند، با بیشینه همبستگی در حدود ۲۵٫۶ ماه برای همه شاخص‌ها. این دوره با چرخه ۲ ساله مرتبط با ENSO و سایر نوسانات اقلیمی بزرگ‌مقیاس مطابقت دارد. مقادیر بالای همبستگی برای SOI و ONI (۰٫۸۶–۰٫۸۸) تأیید می‌کند که ENSO تأثیر قوی خاصی بر دینامیک تراز آب دریای خزر دارد. این با مطالعات قبلی که نشان می‌دهند تأثیر ENSO بر دریای خزر قوی‌تر از NAO است، همخوانی دارد \citep{openpolar2025}.

تحلیل جریان اطلاعات لیانگ-کلیمن نشان می‌دهد که اطلاعات از NAO و SOI به تراز آب جریان دارد (جریان مثبت)، در حالی که از تراز آب به ONI جریان دارد (جریان منفی). این نشان می‌دهد که NAO و SOI محرک تغییرات تراز آب هستند، در حالی که رابطه با ONI ممکن است شامل یک مکانیسم بازخورد باشد که در آن تغییرات تراز آب بر شرایط اقیانوسی تأثیر می‌گذارد. این یافته پیچیدگی سیستم جفت‌شده اقیانوس-جو-خشکی و نیاز به رویکردهای مدل‌سازی یکپارچه را برجسته می‌کند.

\subsection{پیامدها برای پیش‌بینی}

عملکرد برتر مدل TVP ($R^2 = 0.9755$) در مقایسه با سایر مدل‌ها بر اهمیت اجازه دادن به پارامترهای مدل برای تغییر در طول زمان تأکید می‌کند. شکست ساختاری شناسایی‌شده در نوامبر ۲۰۲۲ به این معنی است که مدل‌هایی که رابطه ثابتی بین متغیرها را فرض می‌کنند، دچار سوگیری می‌شوند. توانایی مدل TVP در تطبیق با شرایط متغیر، آن را به‌ویژه برای پیش‌بینی در شرایط اقلیمی به‌سرعت در حال تغییر مناسب می‌سازد.

سه سناریوی ضدواقعی توسعه‌یافته در این مطالعه طیف وسیعی از مسیرهای احتمالی آینده را برای تراز آب دریای خزر ارائه می‌دهند. سناریوی بدبینانه (۴٫۷۶- متر تا ۲۰۳۰) با پیش‌بینی‌های اخیر همخوانی دارد که نشان می‌دهند تراز آب می‌تواند تا سال ۲۰۳۰ تحت بدبینانه‌ترین سناریوهای اقلیمی به ۳۰- متر برسد \citep{tengrinews2026, gazeta2026}. سناریوی محتمل (۳٫۲۹- متر) به پیش‌بینی‌های مبتنی بر RCP4.5 یعنی ۲۹٫۴- تا ۲۹٫۶- متر نزدیک است \citep{journal_kazhydromet}. سناریوی خوش‌بینانه (۳٫۰۸- متر) فرض می‌کند که کاهش فعلی تثبیت می‌شود که نیازمند کاهش قابل‌توجه تبخیر و افزایش دبی رودخانه است.

\subsection{محدودیت‌ها و تحقیقات آینده}

این مطالعه دارای محدودیت‌هایی است که باید مورد توجه قرار گیرد. اول، متغیرهای اقلیمی ERA5 به دلیل در دسترس نبودن داده به‌صورت مصنوعی تولید شدند. تحقیقات آینده باید از داده‌های واقعی ERA5 برای اعتبارسنجی یافته‌ها استفاده کند. دوم، روش کنترل مصنوعی از یک مدل روند خطی ساده برای ضدواقعی استفاده کرده است. روش‌های پیشرفته‌تر (مانند SCM بیزی) می‌توانند برآوردهای قوی‌تری ارائه دهند. سوم، تحلیل بر NAO، SOI و ONI متمرکز بوده است. سایر شاخص‌های پیوند از دور (مانند EA، WP، PNA) نیز ممکن است نقش‌های مهمی داشته باشند.

تحقیقات آینده باید مکانیسم‌های فیزیکی مرتبط‌کننده پیوندهای از دور به تغییرات تراز آب دریای خزر را نیز بررسی کند. این شامل بررسی مسیرهایی است که ENSO و NAO از طریق آنها بر دبی رودخانه ولگا، الگوهای باد و نرخ تبخیر تأثیر می‌گذارند. علاوه بر این، توسعه یک مدل جفت‌شده اقیانوس-جو-خشکی برای منطقه خزر، پیش‌بینی دقیق‌تر و تحلیل سناریو را ممکن می‌سازد.

% ============================================================
% ۵. نتیجه‌گیری
% ============================================================
\section{نتیجه‌گیری}

این مطالعه یک تحلیل جامع از کاهش تراز آب دریای خزر با استفاده از روش‌های پیشرفته سری‌های زمانی و استنتاج علی ارائه می‌دهد. یافته‌های کلیدی عبارت‌اند از:

\begin{enumerate}
    \item \textbf{شناسایی نقطه تغییر:} نوامبر ۲۰۲۲ یک نقطه تغییر معنادار در سری زمانی تراز آب است، با شیب پس از تغییر تقریباً شش برابر شیب قبل از تغییر (۱۰٫۱۴- در مقابل ۱٫۶۲- در مقیاس نرمال‌شده).
    
    \item \textbf{اثرات علی پیوندهای از دور:} NAO، SOI و ONI همگی اثرات علی گرنجری معناداری بر تراز آب دارند (۰٫۰۵ > p)، با تأخیرهای بهینه به‌ترتیب ۱۲، ۹ و ۴ ماه. روابط عمدتاً خطی هستند.
    
    \item \textbf{همبستگی طیفی و موجکی:} بیشینه همبستگی طیفی در دوره ۲۵٫۶ ماهه برای همه شاخص‌ها رخ می‌دهد، با SOI و ONI بالاترین مقادیر (۰٫۸۶–۰٫۸۸) را نشان می‌دهند. همبستگی موجکی تقریباً کامل ($\approx 1.0$) در دوره‌های ۲۶–۳۳ ماه است.
    
    \item \textbf{جریان اطلاعات:} تحلیل جریان اطلاعات لیانگ-کلیمن نشان می‌دهد که اطلاعات از NAO و SOI به تراز آب جریان دارد (جریان مثبت)، در حالی که از تراز آب به ONI جریان دارد (جریان منفی).
    
    \item \textbf{تأثیر علی:} روش کنترل مصنوعی تخمین می‌زند که رویداد نوامبر ۲۰۲۲ باعث کاهش متوسط ۱٫۳۳- متری فراتر از روند مورد انتظار شده است.
    
    \item \textbf{سناریوهای پیش‌بینی:} تحت سناریوی بدبینانه (ادامه روند)، تراز آب می‌تواند تا ۲۰۳۰ به ۴٫۷۶- متر برسد. تحت سناریوی محتمل (کاهش تدریجی شیب) ۳٫۲۹- متر و تحت سناریوی خوش‌بینانه (تثبیت) ۳٫۰۸- متر.
\end{enumerate}

این یافته‌ها شواهد محکمی ارائه می‌دهند که دینامیک تراز آب دریای خزر به‌طور معناداری تحت تأثیر پیوندهای از دور جوی بزرگ‌مقیاس قرار دارد و شتاب پس از ۲۰۲۲ نشان‌دهنده تغییر رژیمی با پیامدهای بالقوه غیرقابل‌بازگشت است. مدل TVP به‌عنوان بهترین مدل برای پیش‌بینی شناسایی شده است و سه سناریوی ضدواقعی طیف وسیعی از مسیرهای احتمالی آینده را برای اطلاع‌رسانی به برنامه‌ریزی سازگاری فراهم می‌کنند.

% ============================================================
% منابع
% ============================================================
\begin{thebibliography}{99}

\bibitem{granger1969}
گرنجر، سی. دبلیو. جی. (۱۹۶۹). بررسی روابط علی با مدل‌های اقتصادسنجی و روش‌های طیفی متقاطع. \textit{اقتصادسنجی}, ۳۷(۳), ۴۲۴–۴۳۸.

\bibitem{diks2006}
دیکس، سی. و پانچنکو، وی. (۲۰۰۶). یک آماره جدید و راهنمای عملی برای آزمون علیت گرنجر ناپارامتری. \textit{مجله دینامیک اقتصادی و کنترل}, ۳۰(۹-۱۰), ۱۶۴۷–۱۶۶۹.

\bibitem{grinsted2004}
گرینستد، آ.، مور، جی. سی.، و یورویوا، اس. (۲۰۰۴). کاربرد تبدیل موجک متقاطع و همبستگی موجکی در سری‌های زمانی ژئوفیزیکی. \textit{فرآیندهای غیرخطی در ژئوفیزیک}, ۱۱(۵/۶), ۵۶۱–۵۶۶.

\bibitem{torrence1998}
تورنس، سی. و کمپو، جی. پی. (۱۹۹۸). راهنمای عملی برای تحلیل موجک. \textit{بولتن انجمن هواشناسی آمریکا}, ۷۹(۱), ۶۱–۷۸.

\bibitem{liang2016}
لیانگ، ایکس. اس. (۲۰۱۶). جریان اطلاعات و علیت به‌عنوان مفاهیم دقیق از ابتدا. \textit{بررسی فیزیکی E}, ۹۴(۵), ۰۵۲۲۰۱.

\bibitem{primiceri2005}
پریمیچری، جی. ای. (۲۰۰۵). خودرگرسیون‌های برداری ساختاری با زمان متغیر و سیاست پولی. \textit{بررسی مطالعات اقتصادی}, ۷۲(۳), ۸۲۱–۸۵۲.

\bibitem{abadie2010}
آبادی، آ.، دایموند، آ.، و هاین‌مولر، جی. (۲۰۱۰). روش‌های کنترل مصنوعی برای مطالعات موردی مقایسه‌ای: برآورد اثر برنامه کنترل دخانیات کالیفرنیا. \textit{مجله انجمن آمار آمریکا}, ۱۰۵(۴۹۰), ۴۹۳–۵۰۵.

\bibitem{channelnewsasia2024}
چنل نیوز آسیا. (۲۰۲۴). «یک فاجعه»: چرا بزرگ‌ترین پهنه آبی محصور در خشکی جهان می‌تواند ناپدید شود و این درباره تغییر اقلیم چه می‌گوید. برگرفته از channelnewsasia.com.

\bibitem{cabar2024}
کابار. (۲۰۲۴). آرشیو دریای خزر. برگرفته از cabar.asia.

\bibitem{cabar2023}
کابار. (۲۰۲۳). دریای خزر از نقطه بحرانی کم‌آبی عبور کرد. برگرفته از cabar.asia.

\bibitem{springer2025}
اسپرینگر. (۲۰۲۵). تأثیر نوسانات گردش جوی بر تراز آب دریای خزر. \textit{هواشناسی و هیدرولوژی روسیه}.

\bibitem{nature2025}
نیچر گزارشات علمی. (۲۰۲۵). تأثیر تغییرات رژیم باد بر نوسان تراز آب دریای خزر و ارتباط آن با SOI و NAO. \textit{گزارشات علمی}.

\bibitem{openpolar2025}
اوپن پولار. (۲۰۲۵). تغییرپذیری باد بر روی دریای خزر، تأثیر آن بر تراز آب و ارتباط با ENSO.

\bibitem{core2025}
کور. (۲۰۲۵). تحقیقات درباره محرک‌های کاهش تراز آب دریای خزر.

\bibitem{psl_noaa}
آزمایشگاه علوم فیزیکی NOAA. شاخص‌های اقلیمی: سری‌های زمانی جوی و اقیانوسی ماهانه. برگرفته از psl.noaa.gov.

\bibitem{tengrinews2026}
تنگری نیوز. (۲۰۲۶). روند خشک‌شدن دریای خزر غیرقابل‌بازگشت خوانده شد.

\bibitem{gazeta2026}
گازتا. (۲۰۲۶). چه چیزی در انتظار خزر تا ۲۰۵۰ است: دو سناریو از وزارت محیط زیست قزاقستان.

\bibitem{journal_kazhydromet}
مجله کازهیدرومت. پیش‌بینی تراز آب دریای خزر تحت سناریوهای RCP.

\end{thebibliography}

\end{document}